\documentclass[11pt]{article}
\usepackage[utf8]{inputenc}
\usepackage[T1]{fontenc}
\setlength{\parindent}{0pt}

\date{\today}
\begin{document}
Hello Giulia,\\
Just a small report on what we’ve been doing so far.\\

\paragraph*{Recap\\}

As discussed with you, our idea was to extend the techniques of Mode Connectivity presented by Garipov.
In their paper,  they focus on simpler curves for mode connectivity - e.g. Bézier curves - with an application to fast ensembling. We aim to study more complex curves, with an application to model merging.\\

We consider the model 1 and 2 - with the same architecture - with respective weights $W_1$ and $W_2$. Both those weights correspond to a mode of the multitask loss.
To connect those modes, we define the following curve
\begin{equation}
    \psi_\theta (t, W_1, W_2) = (1-t) \cdot W_1 + t \cdot W_2 + (1-t)*t \cdot f_\theta(W_1, W_2, t)
\end{equation}
where $f_\theta$ is an MLP parameterized by $\theta$. For now, $\psi$ takes value in the weight space. We refer to the outputs of this function as the interpolated weights.\\

The MLP $f$ is trained as follow: we sample $n$ points $t_i$ uniformly between 0 and 1, compute $l_i$ the loss corresponding to $\psi_\theta (t_i, W_1, W_2)$, and finally optimize $\theta$ with regard to the loss $\sum_{i=1}^{n} l_i$.

\paragraph*{Progress}
\begin{itemize}
  \item We've implemented the curve $\psi$ as well as the above training protocol. The use of the functorch library was practical.
  \item We've started experimenting with Weights and Biases to try to organize our experiments.
  \item We've ran some experiments already and it's performing better than Bézier curves for class incremental learning. We'll share more data soon.
\end{itemize}

\paragraph*{To come}
\begin{itemize}
  \item More experiments and printing graphs of the result to get a better understanding of the shapes of our curves. 
  \item Modifying our optimizer: Since we're sampling $t$ between 0 and 1, the loss can be unstable.
  \item The curve is probably highly non-smooth. Maybe experience with some Smoothness Regularizers that we would add to the loss function.
\end{itemize}

Not easy to progress too fast on this project since we both have a lot on our plate this semester but we still wanted to share the little progress we made. Any feedback is welcomed :)
\end{document}


